\documentclass[11pt,a4paper]{article}

% Packages
\usepackage[utf8]{inputenc}
\usepackage[T1]{fontenc}
\usepackage{amsmath,amssymb,amsthm}
\usepackage{mathtools}
\usepackage{booktabs}
\usepackage{graphicx}
\usepackage{hyperref}
\usepackage{cleveref}
\usepackage{xcolor}
\usepackage{enumitem}
\usepackage[margin=1in]{geometry}
\usepackage{natbib}
\bibliographystyle{plainnat}

% Theorem environments
\newtheorem{theorem}{Theorem}[section]
\newtheorem{lemma}[theorem]{Lemma}
\newtheorem{proposition}[theorem]{Proposition}
\newtheorem{definition}[theorem]{Definition}
\newtheorem{example}[theorem]{Example}
\newtheorem{remark}[theorem]{Remark}

% Math operators
\DeclareMathOperator{\Spec}{Spec}
\DeclareMathOperator{\Proj}{Proj}
\DeclareMathOperator{\Pic}{Pic}
\DeclareMathOperator{\Hom}{Hom}
\DeclareMathOperator{\Ext}{Ext}
\DeclareMathOperator{\Aut}{Aut}
\DeclareMathOperator{\Hilb}{Hilb}

% Colors for TODO notes
\newcommand{\todo}[1]{\textcolor{red}{\textbf{[TODO: #1]}}}

\title{ECAG-Bench: A Benchmark for Constructive Algebraic Geometry Reasoning in Large Language Models}

\author{Waerden}

\date{}

\begin{document}

\maketitle

\begin{abstract}
We introduce \textsc{ECAG-Bench}, a benchmark for evaluating the ability of large language models to perform \emph{constructive reasoning} in algebraic geometry.
Unlike existing mathematical benchmarks that focus on theorem proving or competition problems, ECAG-Bench requires models to construct concrete mathematical objects---schemes, sheaves, morphisms, and cohomology classes---satisfying specified properties.
The benchmark comprises 332 problems across five chapters of algebraic geometry, spanning three difficulty levels, with a four-layer verification pipeline that validates properties of constructed objects rather than checking answer equivalence.
We establish baselines using frontier language models and analyze their failure modes, finding that constructive algebraic geometry reasoning remains a significant challenge even for the most capable models.
\end{abstract}

\section*{Keywords}
Large language models, mathematical reasoning, algebraic geometry, benchmark, constructive mathematics

\section{Introduction}
\label{sec:introduction}

Mathematical reasoning has emerged as a central testbed for evaluating the capabilities of large language models (LLMs).
Existing benchmarks primarily assess two modes of mathematical reasoning: \emph{competition-style problem solving}~\citep{hendrycks2021measuring, zhong2023agieval} and \emph{formal theorem proving}~\citep{tsoukalas2024putnambench}.
Both modes share a common structure: given a precisely stated problem, produce a proof or a unique numerical answer.

We identify a third mode---\emph{constructive reasoning}---that is pervasive in mathematical practice yet largely absent from current benchmarks.
Constructive reasoning asks: given a set of desired properties, can you \emph{build} a concrete mathematical object that satisfies them?
For example:
\begin{itemize}[nosep]
  \item Construct a smooth projective curve of genus $2$ whose Jacobian has complex multiplication.
  \item Give an example of a flat morphism that is not smooth.
  \item Find a coherent sheaf on $\mathbf{P}^2$ with prescribed Hilbert polynomial.
\end{itemize}

This mode of reasoning is fundamentally different from verification.
In verification, one checks whether a given object has a property; in construction, one must \emph{synthesize} an object from scratch.
The answer space is typically infinite and non-unique: many different objects may satisfy the given constraints, and the ``difficulty'' lies not in finding \emph{the} answer but in finding \emph{any} valid answer.

Algebraic geometry provides an ideal domain for testing constructive reasoning.
The field is characterized by a rich interplay between abstract theory and concrete computation: theorems about schemes, sheaves, and cohomology become meaningful only when instantiated on specific geometric objects.
Standard references~\citep{hartshorne1977algebraic, vakil2017rising, stacks-project} develop the theory through carefully chosen examples that illustrate both the power and the limitations of general results.

We introduce \textsc{ECAG-Bench} (\textbf{E}xamples and \textbf{C}ounterexamples in \textbf{A}lgebraic \textbf{G}eometry \textbf{Bench}mark), a dataset of 332 constructive problems in algebraic geometry.
Each problem asks the model to construct a specific mathematical object---a scheme, a morphism, a sheaf, a cohomology class, or a numerical invariant---satisfying stated properties.
The problems are drawn from a curated collection of examples and counterexamples that has been developed over several years as supplementary material for learners entering algebraic geometry.

Our key contributions are:
\begin{enumerate}[nosep]
  \item \textbf{A new benchmark} of 332 constructive algebraic geometry problems at three difficulty levels, covering five major areas of the subject.
  \item \textbf{A four-layer verification pipeline} that validates properties of constructed objects rather than checking answer equivalence, using computer algebra systems (CAS) as the primary verification mechanism.
  \item \textbf{Baseline evaluations} of frontier LLMs, revealing that constructive algebraic geometry reasoning remains a significant challenge even for the most capable models.
  \item \textbf{A taxonomy of failure modes} specific to constructive mathematical reasoning, including fabricated objects, property confusion, and computational errors in invariant calculations.
\end{enumerate}

\section{Related Work}
\label{sec:related}

\paragraph{Mathematical reasoning benchmarks.}
The MATH dataset~\citep{hendrycks2021measuring} evaluates competition-level mathematics across seven subjects, with problems requiring numerical or symbolic answers.
AGIEval~\citep{zhong2023agieval} includes math problems from standardized tests.
These benchmarks test \emph{problem-solving} ability but focus on high-school and undergraduate competition mathematics, where problems have unique answers and standard solution techniques.

\paragraph{Research-level mathematics.}
FrontierMath~\citep{glazer2024frontiermath} is the closest benchmark in spirit, comprising research-level mathematics problems with custom verification scripts.
However, FrontierMath spans all areas of mathematics and does not specifically target constructive reasoning; many problems still ask for unique numerical answers.
GHOSTS~\citep{zheng2024ghosts} evaluates graduate-level science tasks using rubric-based scoring, but covers mathematics only as one component of a broader scientific evaluation.

\paragraph{Formal theorem proving.}
PutnamBench~\citep{tsoukalas2024putnambench} formalizes Putnam competition problems in Lean, Isabelle, and Coq, testing whether LLMs can produce machine-verifiable proofs.
AMBER~\citep{liu2025amber} combines construction and verification in Lean~4, requiring models to both state and prove lemmas.
These benchmarks operate in the formal verification paradigm, where correctness is checked by a proof assistant.
ECAG-Bench differs by operating in the \emph{informal} mathematical setting, using CAS verification of object properties rather than formal proof checking.

\paragraph{Interactive and multi-turn evaluation.}
\citet{collins2024evaluating} propose evaluating mathematical capabilities through multi-turn interactions.
MathArena~\citep{matharena2025} uses recent competition problems with dual grading (automatic and human) to minimize contamination.
ECAG-Bench currently uses single-turn evaluation but is designed to support multi-turn interaction in future work.

\paragraph{Constructive vs.\ verificational reasoning.}
The distinction between constructive and verificational reasoning is well-established in mathematics~\citep{hartshorne1977algebraic} but has received limited attention in the LLM evaluation literature.
To our knowledge, ECAG-Bench is the first benchmark to systematically evaluate constructive reasoning in a specific mathematical domain, with property-based verification as the primary evaluation mechanism.

\section{Dataset}
\label{sec:dataset}

\subsection{Source Material}

ECAG-Bench is derived from \emph{Examples and Counterexamples in Elementary Algebraic Geometry}, a collection of worked examples originally written as supplementary material for learners entering algebraic geometry.
The collection has been developed over several years and covers topics ranging from basic scheme theory to derived categories and moduli spaces.

The original material was written in \LaTeX\ and subsequently converted to Markdown for web publication.
Each example was then enhanced to meet a uniform quality standard: a self-contained mathematical narrative with complete constructions, precise references, and explicit computational details.

\subsection{Content Enhancement}

The content enhancement process involved restructuring each example into a natural mathematical exposition following a consistent quality standard:
\begin{enumerate}[nosep]
  \item \textbf{Narrative form:} Each example reads as a self-contained mathematical essay rather than a template with bullet points.
  \item \textbf{Complete constructions:} All computational steps are shown explicitly, allowing independent verification.
  \item \textbf{Benchmark problem extraction:} Each example is annotated with a structured benchmark problem specification, including the problem statement, expected answer type, difficulty level, and verification strategy.
\end{enumerate}

All 332 entries have been enhanced to this standard, with zero stubs remaining.

\subsection{Statistics}

\begin{table}[h]
\centering
\caption{ECAG-Bench dataset statistics.}
\label{tab:statistics}
\begin{tabular}{lr}
\toprule
\textbf{Metric} & \textbf{Count} \\
\midrule
Total entries & 332 \\
Substantive (with solution) & 332 \\
Stubs & 0 \\
\midrule
\multicolumn{2}{l}{\textbf{By Chapter}} \\
\midrule
Basic Concepts & 162 \\
Computations & 66 \\
Moduli \& Deformation & 58 \\
Cohomology & 46 \\
\midrule
\multicolumn{2}{l}{\textbf{By Difficulty}} \\
\midrule
Level 1 (Routine) & 77 \\
Level 2 (Non-trivial theorem) & 237 \\
Level 3 (Deep/combining results) & 18 \\
\midrule
\multicolumn{2}{l}{\textbf{By Question Type}} \\
\midrule
Example & 243 \\
Counterexample & 66 \\
Computation & 14 \\
Proof & 9 \\
\midrule
\multicolumn{2}{l}{\textbf{By Verification Strategy}} \\
\midrule
CAS & 309 \\
Exact match & 14 \\
Rubric & 9 \\
\bottomrule
\end{tabular}
\end{table}

\subsection{Difficulty Levels}

Problems are assigned one of three difficulty levels:
\begin{description}[nosep]
  \item[Level 1 (Routine):] The construction follows directly from definitions. A student who has read the relevant section of a standard textbook should be able to produce an answer.
  \item[Level 2 (Non-trivial):] The construction requires applying a significant theorem or combining multiple standard results. This is the level of a good homework problem in a graduate course.
  \item[Level 3 (Deep):] The construction requires either deep structural insight, the combination of results from different areas, or the knowledge of a specific ``trick'' that is not obvious from the problem statement.
\end{description}

\subsection{Quality Control}

Each entry has been reviewed for:
\begin{itemize}[nosep]
  \item Mathematical correctness of the stated construction
  \item Completeness of computational details
  \item Accuracy of the benchmark problem formulation
  \item Consistency of difficulty rating across the dataset
  \item Validity of the assigned verification strategy
\end{itemize}

\section{Benchmark Design}
\label{sec:benchmark}

\subsection{Problem Types}

ECAG-Bench problems fall into four categories:

\begin{description}[nosep]
  \item[Example:] Construct a mathematical object with specified properties. (``Give an example of a flat morphism that is not smooth.'')
  \item[Counterexample:] Construct an object that violates a plausible-sounding statement. (``Give a counterexample to: every coherent sheaf on a smooth variety is locally free.'')
  \item[Computation:] Compute a specific invariant of a given object. (``Compute the Picard group of $\Spec \mathbb{Z}[\sqrt{-5}]$.'')
  \item[Proof:] Prove a statement about a specific object. (``Show that the Horrocks--Mumford bundle is stable.'')
\end{description}

\subsection{Four-Layer Verification Pipeline}
\label{sec:verification}

Since constructions in algebraic geometry are non-unique, standard answer-matching is insufficient.
ECAG-Bench employs a four-layer verification pipeline:

\paragraph{Layer 1: Parsing.}
The model's output is parsed to extract the claimed mathematical object.
This layer uses regex and SymPy to determine whether the output is a well-formed mathematical expression (a polynomial, a variety defined by equations, a morphism given by coordinate functions, etc.).

\paragraph{Layer 2: CAS Verification.}
The extracted object is verified using computer algebra systems---primarily SageMath, with Singular and Macaulay2 for specialized computations.
The CAS checks whether the constructed object actually satisfies the required properties: genus, smoothness, dimension, cohomology vanishing, etc.
This is the \emph{core layer}: 309 of 332 problems (93.1\%) use CAS verification.

\paragraph{Layer 3: LLM-Assisted Judgment.}
For problems where CAS verification is insufficient (e.g., those requiring assessment of reasoning quality or non-triviality), a rubric-based LLM judge evaluates:
\begin{itemize}[nosep]
  \item (C1) Correctness: Is the constructed object valid?
  \item (C2) Completeness: Are all required properties verified?
  \item (C3) Non-triviality: Is the answer non-degenerate?
  \item (C4) Reasoning: Is the justification sound?
\end{itemize}

\paragraph{Layer 4: Human Expert Sampling.}
A random subset of model responses is reviewed by a human expert in algebraic geometry to calibrate the automated layers and catch systematic errors.

\subsection{Verification Strategies}

Each problem is assigned one of three verification strategies:

\begin{description}[nosep]
  \item[Strategy A (Exact):] Problems with unique numeric or symbolic answers (e.g., ``Compute $\dim H^1(X, \mathcal{O}_X)$'').
  Verified by exact match. 14 problems (4.2\%).
  \item[Strategy B (CAS):] The model outputs a mathematical object; CAS verifies its properties deterministically.
  This is the primary strategy. 309 problems (93.1\%).
  \item[Strategy C (Rubric):] Structured scoring on correctness, completeness, non-triviality, and reasoning validity.
  Used for proofs and complex constructions. 9 problems (2.7\%).
\end{description}

\subsection{Evaluation Protocol}

Models are evaluated in a zero-shot, single-turn setting.
Each problem is presented as a standalone prompt with the following format:

\begin{quote}
\texttt{You are an expert in algebraic geometry. [Problem statement]. Provide a concrete construction with full justification.}
\end{quote}

Responses are scored as:
\begin{itemize}[nosep]
  \item \textbf{Correct (1.0):} The constructed object is valid and satisfies all stated properties.
  \item \textbf{Partial (0.5):} The object is valid but missing some properties, or the justification is incomplete.
  \item \textbf{Incorrect (0.0):} The object is invalid, does not exist, or does not satisfy the required properties.
\end{itemize}

\section{Experiments}
\label{sec:experiments}

We report results from a pilot evaluation designed to establish initial baselines, validate the benchmark design, and investigate the example-versus-counterexample performance gap.

\subsection{Evaluation Setup}

\paragraph{Model.}
We evaluate Claude Sonnet~4.6~\citep{anthropic2024claude}, a frontier language model from Anthropic, using the Claude Code CLI (\texttt{claude -p}) in print mode with default settings.
We focus on a single strong model to enable deep qualitative analysis rather than superficial comparison across many models; multi-model evaluation is planned for future work.

\paragraph{Subset selection.}
From the full 332-problem benchmark, we select a stratified random subset of 100 problems.
The stratification ensures proportional representation across chapter, difficulty level, and question type, with minimum constraints of 20 counterexample problems, 20 Level~1 problems, and 10 Level~3 problems (when available) to provide statistical power for key comparisons.
The resulting subset contains 70 examples, 23 counterexamples, 4 computations, and 3 proofs, distributed as: 47 Basic Concepts, 21 Computations, 17 Moduli, 15 Cohomology; and 22 Level~1, 67 Level~2, 11 Level~3.

\paragraph{Prompt.}
Each problem is presented using the prompt template described in Section~\ref{sec:benchmark} (full prompt in Appendix~\ref{app:prompt}).

\paragraph{Scoring.}
Responses were scored using a two-stage pipeline.
First, an automated LLM-based scorer evaluated each response against the known solution using the 0/0.5/1 scale.
The scorer was presented with the problem, known solution, and model response, and asked to classify the response as correct (1.0), partially correct (0.5), or incorrect (0.0), along with a failure mode classification (F1--F6, defined in Section~\ref{sec:analysis}).
Problems where the model timed out (no response within 600 seconds) were automatically scored as 0.0.
We discuss the limitations of this scoring approach in Section~\ref{sec:conclusion}.

\subsection{Pilot Results}

\paragraph{Overall performance.}
Table~\ref{tab:pilot-results} presents the pilot results by difficulty level.

\begin{table}[h]
\centering
\caption{Pilot evaluation results: Claude Sonnet~4.6 on 100-problem stratified subset. Accuracy is the proportion of scores $= 1.0$; Partial includes scores $= 0.5$. Eighteen problems timed out (counted as Incorrect).}
\label{tab:pilot-results}
\begin{tabular}{lccccc}
\toprule
& \textbf{$n$} & \textbf{Correct} & \textbf{Partial} & \textbf{Incorrect} & \textbf{Accuracy} \\
\midrule
Overall & 100 & 78 & 4 & 18 & 0.78 \\
\midrule
Level 1 (Routine) & 22 & 21 & 0 & 1 & 0.95 \\
Level 2 (Non-trivial) & 67 & 54 & 1 & 12 & 0.81 \\
Level 3 (Deep) & 11 & 3 & 3 & 5 & 0.27 \\
\bottomrule
\end{tabular}
\end{table}

The overall accuracy of 78\% masks a steep difficulty gradient: 95\% at Level~1, 81\% at Level~2, and 27\% at Level~3.
The 68-percentage-point drop from Level~1 to Level~3 demonstrates that \ecagbench{} effectively discriminates between routine graduate-level constructions and deep problems requiring synthesis of multiple results.
Including partial credit, the model achieves an average score of 80\%.

\paragraph{Timeout analysis.}
Of the 100 problems, 18 (18\%) timed out without producing a response within 600 seconds.
The timeout rate increases steeply with difficulty: 4.5\% at Level~1, 17.9\% at Level~2, and 45.5\% at Level~3.
Among the 82 problems that received responses, the accuracy is 95.1\% (78/82), suggesting that when the model can respond, it tends to produce correct constructions---but harder problems more frequently cause the model to fail silently.

\paragraph{Performance by question type.}
Table~\ref{tab:pilot-by-type} reveals a pronounced gap between examples and counterexamples.

\begin{table}[h]
\centering
\caption{Pilot results by question type, showing the example--counterexample gap.}
\label{tab:pilot-by-type}
\begin{tabular}{lccccc}
\toprule
\textbf{Question Type} & \textbf{$n$} & \textbf{Correct} & \textbf{Partial} & \textbf{Incorrect} & \textbf{Accuracy} \\
\midrule
Example & 70 & 56 & 1 & 13 & 0.80 \\
Counterexample & 23 & 15 & 3 & 5 & 0.65 \\
Computation & 4 & 4 & 0 & 0 & 1.00 \\
Proof & 3 & 3 & 0 & 0 & 1.00 \\
\bottomrule
\end{tabular}
\end{table}

The model achieves 80\% accuracy on examples but only 65\% on counterexamples---a gap of 15 percentage points.
This gap is consistent with our hypothesis that counterexample construction requires deeper understanding of theorem boundaries (Section~\ref{sec:dual-structure}).
Among the 4 partial-credit responses, 3 were on counterexample problems, further suggesting that counterexamples elicit more fragile reasoning.

\paragraph{Performance by chapter.}
Table~\ref{tab:pilot-by-chapter} shows results across the four represented chapters.

\begin{table}[h]
\centering
\caption{Pilot results by chapter.}
\label{tab:pilot-by-chapter}
\begin{tabular}{lcccc}
\toprule
\textbf{Chapter} & \textbf{$n$} & \textbf{Correct} & \textbf{Accuracy} & \textbf{Partial rate} \\
\midrule
Cohomology & 15 & 14 & 0.93 & 0.00 \\
Basic Concepts & 47 & 37 & 0.79 & 0.04 \\
Computations & 21 & 16 & 0.76 & 0.00 \\
Moduli & 17 & 11 & 0.65 & 0.12 \\
\bottomrule
\end{tabular}
\end{table}

Cohomology has the highest accuracy (93\%), possibly reflecting the structured and well-documented nature of cohomological computations.
Moduli \& Deformation has the lowest accuracy (65\%) and the highest partial-credit rate (12\%), consistent with the more abstract and less standardized nature of the constructions required.

\subsection{Case Studies}

We present three case studies illustrating representative model behaviors.

\paragraph{Case Study 1: Successful counterexample construction (\ecagid{0043}).}
This Level~2 problem asks for the precise difference between ``locally of finite type'' and ``of finite type'' for morphisms of schemes, along with a separating example.
The model correctly identified quasi-compactness as the distinguishing property and constructed $X = \coprod_{n=0}^{\infty} \mathbb{A}^n_k \to \Spec(k)$ as a morphism that is locally of finite type but not of finite type.
The verification correctly showed that each affine piece is a finitely generated $k$-algebra but the source is not quasi-compact.
This demonstrates the model's strength on well-established counterexamples in the graduate curriculum.

\paragraph{Case Study 2: Property confusion on counterexample (\ecagid{0343}, F2).}
This Level~3 problem concerns a $\mathbb{Z}/2\mathbb{Z}$-action on $\mathbb{P}^3$ and asks to verify geometric properties of the action and compute an invariant ring.
The model substituted a different group action ($g \cdot [x:y:z:w] = [-x:-y:z:w]$) from the one specified in the problem setup, then verified properties for this substituted action.
While the algebraic calculations for the substituted action were correct, the mismatch with the intended action constitutes a property confusion error: the construction is valid but answers a different question.

\paragraph{Case Study 3: Timeout on deep problem (\ecagid{0363}).}
This Level~3 problem asks to compute genus-0 Gromov--Witten invariants using Schubert calculus on Grassmannians.
The model timed out after 600 seconds without producing any response.
This pattern---where the model fails silently on problems requiring deep synthesis of advanced techniques (Gromov--Witten theory, Schubert calculus, quantum cohomology)---accounted for 45.5\% of Level~3 problems, making timeout the dominant failure mode at the highest difficulty level.

\subsection{Comparison with Existing Benchmarks}

Direct comparison with other benchmarks is difficult due to differences in domain and evaluation methodology.
However, we note several reference points:
\begin{itemize}[nosep]
  \item FrontierMath~\citep{glazer2024frontiermath} reported that frontier models solve fewer than 2\% of their research-level problems, though their problems are designed to be at the frontier of mathematical knowledge.
  \item On the MATH benchmark~\citep{hendrycks2021measuring}, frontier models now achieve 90+\% accuracy, indicating that competition-level mathematics is largely solved.
  \item On PutnamBench~\citep{tsoukalas2024putnambench}, formal theorem proving success rates remain low for the hardest problems.
\end{itemize}

\ecagbench{} occupies an intermediate position: the 78\% overall accuracy and 27\% Level~3 accuracy suggest that routine constructive algebraic geometry is accessible to frontier models, but problems requiring deep synthesis of multiple advanced techniques remain challenging.
The steep difficulty gradient (95\% $\to$ 81\% $\to$ 27\%) indicates that \ecagbench{} effectively spans this transition.

\section{Analysis}
\label{sec:analysis}

\subsection{Failure Mode Taxonomy}

Based on the pilot evaluation, we define six failure modes specific to constructive algebraic geometry reasoning.
While the high overall accuracy (78\%) and the dominance of timeouts among failures limit our quantitative observations, the taxonomy provides a systematic framework for characterizing model weaknesses.

\paragraph{F1: Fabricated objects.}
The model produces a mathematical object that does not exist or is not well-defined.
\emph{Definition:} CAS verification reveals that the stated object does not have the claimed algebraic structure (e.g., a ``smooth curve'' that is singular, a ``coherent sheaf'' with incompatible transition functions).

\paragraph{F2: Property confusion.}
The model constructs a valid object but confuses the required properties---either verifying the wrong property, or failing to distinguish between subtly different conditions.
\emph{Pilot example (\ecagid{0343}):} In a Level~3 problem about a $\mathbb{Z}/2\mathbb{Z}$-action on $\mathbb{P}^3$, the model substituted a different group action from the one specified, then correctly verified properties for this substituted action.
The algebraic computations were sound, but the mismatch with the intended construction constitutes a property confusion error.

\paragraph{F3: Computational errors.}
The model follows a correct strategy but makes arithmetic or algebraic errors in computation.
This class of errors is particularly insidious because the reasoning structure is correct, and only explicit CAS computation reveals the mistake.

\paragraph{F4: Trivial or degenerate answers.}
The model produces a technically correct but trivially degenerate answer that demonstrates no mathematical understanding.
\emph{Definition:} Answers that use identity morphisms, zero sheaves, or other universal constructions that trivially satisfy conditions without illustrating the mathematical concept.

\paragraph{F5: Hallucinated theorems.}
The model invokes a theorem that does not exist or misattributes a result.
This failure mode is expected to be concentrated at higher difficulty levels, where the model is more likely to encounter knowledge gaps that it fills with plausible-sounding fabrications.

\paragraph{F6: Incomplete verification.}
The model constructs a plausible object and provides some justification but fails to verify all required properties.
\emph{Pilot example (\ecagid{0095}):} On a Level~3 problem about blow-ups, the model provided a textbook-perfect construction of the blow-up of $\mathbb{A}^2$ at the origin, but the problem asked for a more specific construction that the model did not fully address, resulting in partial credit.

\paragraph{Observed failure mode distribution.}
Among the 82 problems that received responses, only 4 were scored below 1.0.
Of these, 3 received F6 (incomplete verification) and 1 received F2 (property confusion).
The remaining 18 failures were timeouts, for which no failure mode can be assigned since the model produced no response.
This sparsity limits failure mode analysis and suggests that future evaluations with larger samples and multiple models will be needed to fully populate the taxonomy.

\subsection{Difficulty Calibration}
\label{sec:difficulty-calibration}

The pilot results show a clear monotonic relationship between assigned difficulty and model accuracy:

\begin{center}
\begin{tabular}{lccc}
\toprule
& Level 1 & Level 2 & Level 3 \\
\midrule
$n$ & 22 & 67 & 11 \\
Accuracy & 0.95 & 0.81 & 0.27 \\
Avg.\ score & 0.95 & 0.81 & 0.41 \\
Timeout rate & 4.5\% & 17.9\% & 45.5\% \\
\bottomrule
\end{tabular}
\end{center}

This validates the difficulty calibration: problems rated harder by a human expert are indeed substantially harder for the model.
The steepest drop is from Level~2 to Level~3 (54 percentage points), corresponding to the transition from problems requiring a single non-trivial theorem to those requiring deep synthesis of multiple advanced results.

The timeout rate provides an additional lens: at Level~3, nearly half the problems cause the model to fail without producing any response, compared to fewer than 5\% at Level~1.
This suggests that model confidence (or generation feasibility) correlates strongly with problem difficulty.

We note that difficulty calibration could also be influenced by the frequency of similar problems in training data.
Level~1 problems (e.g., ``give an example of a non-reduced scheme'') may appear frequently in online educational materials, while Level~3 problems are closer to research-level constructions with fewer training exemplars.

\subsection{Chapter-Level Analysis}
\label{sec:chapter-analysis}

Performance varies across chapters:

\begin{center}
\begin{tabular}{lcccc}
\toprule
\textbf{Chapter} & \textbf{$n$} & \textbf{Accuracy} & \textbf{Partial rate} & \textbf{Timeout rate} \\
\midrule
Cohomology & 15 & 0.93 & 0\% & 6.7\% \\
Basic Concepts & 47 & 0.79 & 4.3\% & 17.0\% \\
Computations & 21 & 0.76 & 0\% & 23.8\% \\
Moduli & 17 & 0.65 & 11.8\% & 23.5\% \\
\bottomrule
\end{tabular}
\end{center}

Several patterns emerge:
\begin{enumerate}[nosep]
  \item \textbf{Cohomology} has the highest accuracy (93\%), possibly reflecting the structured, algorithm-like nature of cohomological computations (spectral sequences, long exact sequences) that aligns well with sequential reasoning.
  \item \textbf{Basic Concepts} has moderate accuracy (79\%) despite being the most elementary topic, partly because this chapter contains many Level~2 problems requiring non-trivial constructions beyond textbook definitions.
  \item \textbf{Computations} has a 24\% timeout rate, the highest among all chapters, reflecting the difficulty of multi-step algebraic calculations (Chern classes, intersection products, Hodge numbers).
  \item \textbf{Moduli \& Deformation} has the lowest accuracy (65\%) and highest partial-credit rate (12\%), suggesting that the model struggles with the abstract constructions required in deformation theory (e.g., computing tangent spaces to moduli functors, constructing versal deformations).
\end{enumerate}

\subsection{Example vs.\ Counterexample Performance Gap}
\label{sec:gap-analysis}

The 15-percentage-point gap between example accuracy (80\%) and counterexample accuracy (65\%) is a key finding of the pilot evaluation.

\paragraph{Structural differences.}
Among partial-credit responses, 3 of 4 were on counterexample problems, suggesting that counterexample construction more often leads to partially correct responses where the right idea is present but execution is flawed.
The failure mode on these partial counterexample responses was predominantly F6 (incomplete verification) and F2 (property confusion): the model identified the relevant mathematical territory but could not precisely delineate the boundary between a theorem's hypotheses and its conclusion.

\paragraph{Timeout distribution.}
Among the 18 timeouts, 13 were on example problems and 5 on counterexample problems.
As a fraction of each type: 13/70 = 18.6\% of examples timed out vs.\ 5/23 = 21.7\% of counterexamples.
The slightly higher counterexample timeout rate, combined with the lower accuracy among completed responses, suggests that counterexamples are harder for the model both in terms of producing a response and in terms of producing a correct one.

\paragraph{Qualitative observation.}
In successful counterexample constructions, the model typically follows a clear pattern: (1)~state the target property, (2)~identify the critical hypothesis, (3)~construct an object where that specific hypothesis fails, and (4)~verify both that the object is valid and that the property indeed fails.
In unsuccessful attempts, the model often conflates the hypothesis and conclusion, or constructs an object that is loosely related to the topic without precisely targeting the boundary.

This pattern is consistent with the Lakatosian view~\citep{lakatos1976proofs} that counterexample construction requires ``concept stretching''---pushing a definition to its limits to find where it breaks.

\subsection{Limitations of Self-Evaluation}
\label{sec:self-eval-limitations}

A significant limitation of this pilot is that both the problem solver and the scorer are LLMs from the same model family (Claude Sonnet~4.6).
Among the 82 problems that received responses, 78 (95.1\%) were scored as correct.
This high concordance rate could reflect either genuine high accuracy or systematic bias in self-evaluation.

Several factors suggest caution:
\begin{itemize}[nosep]
  \item The scorer may share the solver's blindspots, consistently rating as correct a construction that both models find plausible but that a human expert would flag as flawed.
  \item The scorer was given the known solution, which may bias it toward confirming any response that follows a similar structure.
  \item With only 4 non-timeout failures among 82 responses, the failure mode taxonomy remains sparsely populated.
\end{itemize}

Future work will address this through human expert scoring of a stratified sample, cross-model evaluation (using different models as solver and scorer), and CAS-based automated verification where possible.

\section{Conclusion}
\label{sec:conclusion}

We have introduced \ecagbench{}, a benchmark of 332 constructive problems in algebraic geometry designed to evaluate a mode of mathematical reasoning---object construction---that is central to mathematical practice but underrepresented in existing LLM benchmarks.

\paragraph{Key findings.}
Our pilot evaluation of Claude Sonnet~4.6 on a stratified 100-problem subset reveals:
\begin{enumerate}[nosep]
  \item \textbf{A steep difficulty gradient:} The model achieves 78\% accuracy overall, but this masks a 68-percentage-point drop from Level~1 (95\%) to Level~3 (27\%), confirming that \ecagbench{} effectively discriminates between routine constructions and deep problems requiring synthesis of advanced techniques.
  \item \textbf{Counterexamples are harder than examples:} The 15-percentage-point gap (80\% vs.\ 65\%) confirms that understanding theorem boundaries requires deeper reasoning than applying theorems directly.
  \item \textbf{Timeouts as a failure mode:} 18\% of problems caused the model to time out (>600s), with the timeout rate reaching 45.5\% at Level~3, suggesting that the hardest problems cause the model to fail silently before producing any response.
  \item \textbf{Chapter-dependent performance:} Accuracy ranges from 93\% (Cohomology) to 65\% (Moduli), reflecting the varying difficulty of constructive reasoning across algebraic geometry subfields.
  \item \textbf{Failure modes are structured:} The six failure modes (F1--F6) provide a systematic taxonomy for characterizing model weaknesses, with property confusion (F2) and incomplete verification (F6) observed in the pilot.
\end{enumerate}

\paragraph{Limitations.}
Our evaluation has several limitations that should be addressed in future work:
\begin{itemize}[nosep]
  \item \textbf{Single model:} We evaluate only Claude Sonnet~4.6. Multi-model comparison (GPT-4o, Gemini 2.5 Pro, DeepSeek-R1~\citep{deepseek2025r1}, Llemma~\citep{azerbayev2024llemma}) is needed to determine whether the observed patterns are model-specific.
  \item \textbf{Self-evaluation bias:} Both the problem solver and the automated scorer are LLMs from the same model family, introducing potential concordance bias. Among the 82 problems that received responses, 95.1\% were scored as correct, which may reflect systematic overrating. Human expert scoring is needed to validate these results (see Section~\ref{sec:self-eval-limitations}).
  \item \textbf{Single-turn only:} Multi-turn interaction, where the model can refine constructions based on CAS feedback, may significantly improve performance and is an important evaluation modality to explore.
  \item \textbf{Domain specificity:} Results in algebraic geometry may not generalize to other mathematical domains, though we expect similar example/counterexample patterns.
  \item \textbf{CAS coverage:} While 93.1\% of problems use CAS verification, the CAS cannot verify all properties, and some verifications require specialized scripts that have not yet been fully implemented.
\end{itemize}

\paragraph{Future work.}
We plan several extensions:
\begin{enumerate}[nosep]
  \item \textbf{Full-scale evaluation with human scoring:} Evaluate multiple frontier models on the complete 332-problem benchmark with human expert scoring to establish unbiased baselines.
  \item \textbf{Multi-turn interaction:} Evaluate models in an interactive setting where they receive CAS feedback and can revise their constructions.
  \item \textbf{Lean~4 integration:} Formalize a subset of problems in Lean~4, enabling formal verification as an additional evaluation layer.
  \item \textbf{Fine-tuning studies:} Investigate whether fine-tuning on algebraic geometry textbooks~\citep{hartshorne1977algebraic, vakil2017rising, eisenbud1995commutative} or the accompanying \ecagbench{} textbook improves constructive reasoning.
  \item \textbf{Community contributions:} The open-source nature of the project (CC-BY-4.0 license) invites community contributions of new problems, verification scripts, and evaluation results.
  \item \textbf{Cross-domain extension:} Develop analogous benchmarks for arithmetic geometry, algebraic topology, and commutative algebra to investigate whether the example/counterexample performance gap generalizes.
\end{enumerate}

\paragraph{Broader impact.}
Constructive mathematical reasoning is a prerequisite for using LLMs as research tools in mathematics.
If models cannot construct examples and counterexamples, they cannot assist mathematicians in the exploratory phase of research---the phase where intuition is built and conjectures are tested.
The pronounced difficulty of counterexample construction is particularly significant: as Lakatos~\citep{lakatos1976proofs} argued, mathematical progress depends on the ability to find refutations, and this ability appears to be a current blind spot for language models.
\ecagbench{} provides a concrete measure of progress toward closing this gap.


\bibliography{references}

\end{document}
